\documentclass[10pt]{article}
\usepackage[T1]{fontenc}
\usepackage[utf8]{inputenc}
\usepackage[french]{babel}
\usepackage{amsmath,amssymb,mathrsfs,mathpazo}
\usepackage{geometry}
\geometry{a4paper,margin=2cm}
\usepackage{graphicx,array,tabularx,tabularray,colortbl,multirow,multicol,booktabs,enumitem,pifont,ulem,fancybox,diagbox,cancel}
\usepackage{pstricks,pstricks-add,pst-plot,pst-tree,pst-node,pst-text,pst-func,pst-math,pst-bspline,pst-3dplot,pst-eucl,pst-solides3d}
\let\pscalculate\undefined
\usepackage{pst-bezier}
\usepackage{tikz,pgfplots,tkz-tab}
\usetikzlibrary{arrows,arrows.meta,patterns,calc,positioning,trees,decorations.markings,shapes.geometric,shapes.misc,shapes.arrows,intersections,decorations.pathmorphing,decorations.pathreplacing,backgrounds,shadows,babel}
\pgfplotsset{compat=1.18}
\usepgfplotslibrary{fillbetween}
\usepackage[np]{numprint}
\usepackage{listings,makecell,pas-tableur,twoopt}
\let\barre\undefined
\usepackage[pstricks]{tablvar}
\definecolor{darkgreen}{RGB}{0,100,0}
\providecommand{\red}{\color{red}}
\providecommand{\blue}{\color{blue}}
\providecommand{\green}{\color{green}}
\providecommand{\black}{\color{black}}
\providecommand{\white}{\color{white}}
\providecommand{\gray}{\color{gray}}
\providecommand{\cyan}{\color{cyan}}
\providecommand{\magenta}{\color{magenta}}
\providecommand{\yellow}{\color{yellow}}
\providecommand{\violet}{\color{violet}}
\providecommand{\orange}{\color{orange}}
\providecommand{\euro}{\texteuro}
\providecommand{\EUR}{\texteuro}
\providecommand{\ieme}{\textsuperscript{e}}
\providecommand{\degres}{\ensuremath{^\circ}}
\providecommand{\degree}{\ensuremath{^\circ}}
\providecommand{\vv}[1]{\overrightarrow{#1}}
\providecommand{\Coord}[1]{#1}
\providecommand{\decosix}{}
\providecommand{\starredbullet}{\ensuremath{\bullet}}
\pagestyle{empty}
\setlength{\parindent}{0pt}
\usepackage{adjustbox}\DeclareUnicodeCharacter{2236}{:}\DeclareUnicodeCharacter{2190}{\ensuremath{\leftarrow}}
\providecommand{\R}{\mathbb{R}}
\providecommand{\N}{\mathbb{N}}
\providecommand{\Z}{\mathbb{Z}}
\providecommand{\Q}{\mathbb{Q}}
\providecommand{\D}{\mathbb{D}}
\providecommand{\C}{\mathbb{C}}
\providecommand{\e}{{\mathrm{e}}}
\providecommand{\vect}[1]{\overrightarrow{#1}}
\providecommand{\vectt}[1]{\overrightarrow{\mathrm{#1}}}
\providecommand{\vvt}[1]{\overrightarrow{\mathrm{#1}}}
\providecommand{\np}[2][]{{#2~#1}}
\providecommand{\vv}[1]{{\overrightarrow{#1}}}
\renewcommand{\d}{{\mathrm{d}}}
\providecommand{\degres}{{^\circ}}
\providecommand{\degre}{{^\circ}}
\providecommand{\blue}{\color{blue}}
\providecommand{\red}{\color{red}}
\providecommand{\enskip}{\quad}
\providecommand{\llbracket}{\lbrack\!\lbrack}
\providecommand{\rrbracket}{\rbrack\!\rbrack}
\providecommand{\euro}{{\text{€}}}
\providecommand{\no}{{\text{n°}}}
\providecommand{\interval}[2]{{[#1;#2]}}
\providecommand{\intervFO}[2]{{[#1;#2[}}
\providecommand{\dotfill}{\ldots}

\providecommand{\ds}{\displaystyle}
\providecommand{\textding}[1]{#1}
\providecommand{\Oij}{\(\mathrm{O};\vec{\imath},\vec{\jmath}\)}
\providecommand{\Oijk}{\(\mathrm{O};\vec{\imath},\vec{\jmath},\vec{k}\)}
\providecommand{\Ouv}{\(\mathrm{O};\vec{u},\vec{v}\)}
\providecommand{\intervFF}[2]{\ensuremath{[#1;#2]}}
\renewcommand{\euro}{\ifmmode\text{\texteuro}\else\texteuro\fi}
\DeclareUnicodeCharacter{00B0}{\ifmmode^{\circ}\else\textdegree\fi}
\DeclareUnicodeCharacter{20AC}{\ifmmode\text{\texteuro}\else\texteuro\fi}
% Énoncés : Éducation nationale. Transcription et corrigés : APMEP, Denis Vergès et contributeurs. Découpage et classement : CoopMaths, Sébastien Lozano, Éric Elter et contributeurs. Adaptation de la présentation : Maths Explorer.
% CC BY-SA : https://coopmaths.fr/www/about/
% https://coopmaths.fr/alea/static/bac/2024/tex/bac_2024_06_sujet2_polynesie_2.tex
\begin{document}
\noindent\textit{APMEP / CoopMaths. CC BY-SA. Présentation adaptée par Maths Explorer.}\par\medskip
\medskip

\textit{Cet exercice est un questionnaire à choix multiples (QCM) qui comprend cinq questions. Les cinq questions sont indépendantes.\\
Pour chacune des questions, \textbf{une seule des quatre réponses est exacte}. Le candidat indiquera sur sa copie le numéro de la question suivi de la lettre correspondant à la réponse exacte.\\
Aucune justification n'est demandée.\\
Une réponse fausse, une réponse multiple ou une absence de réponse ne rapporte, ni n'enlève aucun point.}

\begin{enumerate}
	\item La solution $f$ de l'équation différentielle $y'=-3y+7$ telle que $f(0)=1$ est la fonction définie sur $\R$ par :
	
	\smallskip
	
\begin{center}
\renewcommand\arraystretch{1.9}
\begin{tabularx}{\linewidth}{X X}
\textbf{A. } $f(x) = \e^{-3x}$&\textbf{B. } $f(x) = - \dfrac43 \e^{-3x} + \dfrac73$\\
\textbf{C. } $f(x) = \e^{-3x} + \dfrac73$&\textbf{D.~} $f(x) = - \dfrac{10}{3} \e^{-3x} - \dfrac73$
\end{tabularx}
\end{center}
	\item La courbe d'une fonction $f$ définie sur $[0~;~+\infty[$ est donnée ci-dessous.
	
\begin{center}
\psset{unit=1.4cm}
\begin{pspicture}(-0.5,-0.25)(6.25,4.25)
\psgrid[gridlabels=0pt,subgriddiv=1]
\psaxes[linewidth=1.25pt,labelFontSize=\scriptstyle]{->}(0,0)(0,0)(6.25,4.25)
\psplot[plotpoints=2000,linewidth=1.25pt,linecolor=red]{0}{6.25}{0.00196  x 5 exp mul 0.05764 x  4 exp mul sub 0.62643 x  3 exp mul add 3.00524 x dup mul mul sub 5.43449  x mul add}
\end{pspicture}
\end{center}

Un encadrement de l'intégrale $I = \displaystyle\int_1^5 f(x) \:\text{d}x$ est :

\begin{center}
\begin{tabularx}{\linewidth}{X X}
\textbf{A. } $0 \leqslant I \leqslant 4$&\textbf{B. } $1 \leqslant I \leqslant 5$\\
\textbf{C. } $5 \leqslant I \leqslant 10$&\textbf{D. } $10 \leqslant I \leqslant 15$
\end{tabularx}
\end{center}

\item On considère la fonction $g$ définie sur $\R$ par $g(x) = x^2 \ln \left(x^2 + 4\right) $.

Alors $\displaystyle\int_0^2 g'(x)\:\text{d}x$ vaut, à $10^{-1}$ près :

\begin{center}
\begin{tabularx}{\linewidth}{X X}
\textbf{A. } 4,9&\textbf{B. } $8,3$\\
\textbf{C. } $1,7$&\textbf{D. } $7,5$
\end{tabularx}
\end{center}	
	\smallskip
	
\item Une professeure enseigne la spécialité mathématiques dans une classe de $31$ élèves de terminale.
	
Elle veut former un groupe de 5 élèves. De combien de façons différentes peut-elle former 
un tel groupe de 5 élèves ?

\begin{center}
\begin{tabularx}{\linewidth}{X X}
\textbf{A. } $31^5$&\textbf{B. } $31\times30\times29\times28\times27$\\
\textbf{C. } $31+30+29+28+27$&\textbf{D. } $\dbinom{31}{5}$
\end{tabularx}
\end{center}

	\item La professeure s'intéresse maintenant à l'autre spécialité des 31 élèves de son groupe :
	
\begin{itemize}
	\item 10 élèves ont choisi la spécialité physique-chimie ;
	\item 20 élèves ont choisi la spécialité SES ;
	\item 1 élève a choisi la spécialité LLCE espagnol.
\end{itemize}
	
Elle veut former un groupe de 5 élèves comportant exactement 3 élèves ayant choisi la
	spécialité SES. De combien de façons différentes peut-elle former un tel groupe ?
	
\begin{center}
\renewcommand\arraystretch{1.9}
\begin{tabularx}{\linewidth}{X X}
\textbf{A. } $\displaystyle\binom{20}{3} \times \binom{11}{2}$&\textbf{B. } $\displaystyle\binom{20}{3} + \binom{11}{2}$\\
\textbf{C. } $\displaystyle\binom{20}{3}$&\textbf{D. } $20^3 \times 11^2$
\end{tabularx}
\end{center}
\end{enumerate}

\bigskip

\end{document}
